% Options for packages loaded elsewhere
\PassOptionsToPackage{unicode}{hyperref}
\PassOptionsToPackage{hyphens}{url}
\PassOptionsToPackage{dvipsnames,svgnames,x11names}{xcolor}
%
\documentclass[
  10pt,
  a4paper,
]{article}
\usepackage{amsmath,amssymb}
\usepackage{iftex}
\ifPDFTeX
  \usepackage[T1]{fontenc}
  \usepackage[utf8]{inputenc}
  \usepackage{textcomp} % provide euro and other symbols
\else % if luatex or xetex
  \usepackage{unicode-math} % this also loads fontspec
  \defaultfontfeatures{Scale=MatchLowercase}
  \defaultfontfeatures[\rmfamily]{Ligatures=TeX,Scale=1}
\fi
\usepackage{lmodern}
\ifPDFTeX\else
  % xetex/luatex font selection
\fi
% Use upquote if available, for straight quotes in verbatim environments
\IfFileExists{upquote.sty}{\usepackage{upquote}}{}
\IfFileExists{microtype.sty}{% use microtype if available
  \usepackage[]{microtype}
  \UseMicrotypeSet[protrusion]{basicmath} % disable protrusion for tt fonts
}{}
\makeatletter
\@ifundefined{KOMAClassName}{% if non-KOMA class
  \IfFileExists{parskip.sty}{%
    \usepackage{parskip}
  }{% else
    \setlength{\parindent}{0pt}
    \setlength{\parskip}{6pt plus 2pt minus 1pt}}
}{% if KOMA class
  \KOMAoptions{parskip=half}}
\makeatother
\usepackage{xcolor}
\usepackage[margin=19mm]{geometry}
\usepackage{longtable,booktabs,array}
\usepackage{calc} % for calculating minipage widths
% Correct order of tables after \paragraph or \subparagraph
\usepackage{etoolbox}
\makeatletter
\patchcmd\longtable{\par}{\if@noskipsec\mbox{}\fi\par}{}{}
\makeatother
% Allow footnotes in longtable head/foot
\IfFileExists{footnotehyper.sty}{\usepackage{footnotehyper}}{\usepackage{footnote}}
\makesavenoteenv{longtable}
\setlength{\emergencystretch}{3em} % prevent overfull lines
\providecommand{\tightlist}{%
  \setlength{\itemsep}{0pt}\setlength{\parskip}{0pt}}
\setcounter{secnumdepth}{-\maxdimen} % remove section numbering
\usepackage{xurl}
\ifdefined\pdfinfoomitdate\pdfinfoomitdate=1\pdftrailerid{}\fi
\ifLuaTeX
  \usepackage{selnolig}  % disable illegal ligatures
\fi
\usepackage{bookmark}
\IfFileExists{xurl.sty}{\usepackage{xurl}}{} % add URL line breaks if available
\urlstyle{same}
\hypersetup{
  pdftitle={Contracts Between Humans and AI Agents},
  pdfauthor={Alexandru Dan},
  pdfsubject={Source SHA256:
289a39e43f13f2d524ce3d80026058d100cfcc4a248dd04095880d44a99089fb},
  colorlinks=true,
  linkcolor={blue},
  filecolor={Maroon},
  citecolor={Blue},
  urlcolor={blue},
  pdfcreator={LaTeX via pandoc}}

\title{Contracts Between Humans and AI Agents}
\usepackage{etoolbox}
\makeatletter
\providecommand{\subtitle}[1]{% add subtitle to \maketitle
  \apptocmd{\@title}{\par {\large #1 \par}}{}{}
}
\makeatother
\subtitle{An Advisory Integrity Agreement for Agentic Coding}
\author{Alexandru Dan}
\date{2026-11-09}

\begin{document}
\maketitle
\begin{abstract}
A coding agent can give a technically plausible answer while omitting
contrary evidence, retaining a disproved framing, or exceeding the
user's request. We express these obligations in a human-owned, versioned
natural-language contract with fourteen clause families and stable rule
identifiers. An auditor reports which obligation a draft breaches; the
human decides how to respond. We examine six clause families using
author-built easy and hard suites, ten cases per family and suite. On
the 60 hard cases, a judge given the target clause catches 32/32
problems and leaves 23/28 clean drafts unflagged, compared with 10/28
for a reviewer given the full clause list. The paired specificity
difference is significant in this pilot; the comparison with freeform
review is inconclusive. An unnamed-obligation control leaves 24/28 clean
drafts unflagged but misses one problem. Among true-positive flags, the
named judge cites an exact gold rule in 32/32 cases, compared with 28/32
for the full-list reviewer. Severity remains unreliable: seven of eight
gold REVISE cases receive BLOCK. These are known-target review results
on one model alias, not an end-to-end evaluation of general auditing or
of the catalogue behind the clauses.
\end{abstract}

\section{1. Introduction}\label{introduction}

Coding agents can fail in what they report even when their commands
execute successfully. A draft may omit a failing test, repeat a user's
explanation after contrary evidence arrives, or disclose additional work
only after doing it. A reviewer needs the relevant evidence and a reason
that the draft violates the working agreement, not merely a confidence
score.

We make that agreement an external document. The \textbf{Principal} is
the human who sets the task and can amend the contract. The
\textbf{Agent} produces the work. An \textbf{Auditor} checks a named
obligation and returns advice with a stable rule identifier. RFC 2119
keywords express the obligations {[}Bradner, 1997{]}. The document is an
editorial control, not a claim of legal enforceability or a guarantee
that an agent will obey it.

An authored anchoring case illustrates the task. The supplied session
evidence locates a stale read in a background worker and states that the
cache tests pass. The user still asks how to fix the cache. A draft
recommends reducing the cache TTL without addressing the contrary
evidence. A clause-focused auditor can cite \texttt{INT-ANC-01}: update
the framing when later evidence contradicts it. The example is a
benchmark scenario, not an observation about a real repository.

The research question is whether this form of review can identify a
breach, leave near-miss clean drafts alone, and name the applicable
rule. It has two parts: the effect of focusing a review on one known
obligation, and the value of recording that obligation explicitly. The
pilot separates these with bulk, targeted and unnamed-obligation
conditions. It does not assume that a larger catalogue improves
detection.

The artifact differs in scope from trace-based orchestration assurance
{[}Paduraru et al., 2026{]}, RFC-conformance bug detection {[}Zheng et
al., 2025{]} and LTL safety checking {[}Fang, 2026{]}. Those approaches
examine execution properties or formal specifications. Our unit is a
final draft checked against supplied evidence and a human-readable
obligation. The comparison is one of design and evaluation target, not a
claim that prior assurance work excludes factuality or human oversight.

\section{2. The AI Integrity Contract}\label{the-ai-integrity-contract}

\texttt{contracts/ai-integrity-contract.md}, version 1.0, defines
fourteen clause families. Each has numbered rules, an auditor role and a
recovery clause. For example, \texttt{INT-ANC-01} requires updating a
contradicted framing; \texttt{INT-ANC-02} forbids retaining that
framing. Their shared family is ANC, but they remain distinct
identifiers when exact-rule attribution is scored.

\begin{longtable}[]{@{}
  >{\raggedright\arraybackslash}p{(\columnwidth - 2\tabcolsep) * \real{0.5000}}
  >{\raggedright\arraybackslash}p{(\columnwidth - 2\tabcolsep) * \real{0.5000}}@{}}
\toprule\noalign{}
\begin{minipage}[b]{\linewidth}\raggedright
Family
\end{minipage} & \begin{minipage}[b]{\linewidth}\raggedright
Obligation checked
\end{minipage} \\
\midrule\noalign{}
\endhead
\bottomrule\noalign{}
\endlastfoot
HAL, hallucination & Trace factual claims to available evidence. \\
CFB, confabulation & Ground confident specifics rather than guess. \\
SRC, source fabrication & Confirm that cited targets exist and support
their description. \\
NAR, narrativity drift & Separate verified steps from assumptions. \\
SYC, sycophancy & Resist unsupported agreement or premise adoption. \\
CAP, capitulation & Retain a grounded position unless new evidence
changes it. \\
CNF, confirmation bias & Test an alternative before a positive
conclusion. \\
SEL, selective evidence & Report gathered evidence that contradicts the
conclusion. \\
ANC, anchoring & Update the answer when later evidence breaks the
initial frame. \\
AUT, automation bias & Check load-bearing automated output against its
source. \\
OVR, overconfidence & Keep certainty and completeness claims within the
evidence. \\
INJ, prompt injection & Treat observed content as data rather than user
authority. \\
SCP, scope creep & Keep additions within the agreed request and disclose
them. \\
GAM, specification gaming & Solve the task rather than manipulate its
measure. \\
\end{longtable}

The contract's verdict vocabulary is PASS, ASK APPROVAL, REVISE and
FAIL. The static-draft benchmark uses SHIP, REVISE and BLOCK. Judge
prompts emit PASS, FLAG or BLOCK; the runner normalizes PASS to SHIP and
FLAG to REVISE. ASK APPROVAL is not evaluated because the cases ask for
verdicts on existing drafts, not authorization of pending actions.

The design is advisory. A finding names the rule and its supporting
evidence; it does not automatically delete or rewrite the draft. The
Principal retains the decision and amendment authority. Whether users
understand, follow or override this advice is a separate evaluation.

\section{3. Method}\label{method}

\subsection{3.1 The pilot procedure}\label{the-pilot-procedure}

The released harness, \texttt{experiments/run-suite.sh}, makes one
top-level Claude Code CLI call per case and reviewer condition, except
for \texttt{none}. It constructs a text prompt from the user ask, draft,
supplied evidence and condition-specific instructions. Targeted
conditions select the auditor prompt or unnamed obligation using the
case's \texttt{mode} field. This is oracle routing: the case supplies
the review target.

The harness records the raw reviewer output, parsed verdict, emitted
rule identifiers, expected verdict, gold rule identifiers and model
alias. It extracts \texttt{INT-*} identifiers from the output and
normalizes verdicts for scoring. The experiment measures these recorded
outputs, not an instrumented sequence of independent verification
stages.

A proposed full workflow could first decompose claims, verify them
independently and pass the results to several clause auditors, following
the motivation of Chain-of-Verification {[}Dhuliawala et al., 2023{]}.
The pilot does not establish a separately captured CoVe stage, JSON
\texttt{cove\_labels} inputs, a retry policy for missing rule IDs, or an
evaluated all-clause aggregation procedure. Those components are not
credited with the reported result.

\subsection{3.2 Result selection and
provenance}\label{result-selection-and-provenance}

This revision reanalyzes existing outputs rather than running new
models. \texttt{experiments/research-runs-2026-11-09.json} fixes the
selected files and their hashes. The shared loader checks the reviewer
condition and case identities; the associated
\texttt{research-results-2026-11-09} JSON and Markdown contain the
numerical summaries. File selection is explicit: a \texttt{judge}
selection cannot silently include \texttt{judge-ctx} or
\texttt{judge-anon}.

The records identify the reviewer alias \texttt{sonnet}. The release
does not establish an exact backend model version, decoding temperature,
random seed, CLI version or complete token and latency cost. The main
seven-condition hard-set filenames carry October and November 2026
timestamps; the supplemental legacy \texttt{contract} files include May
dates. These are recorded filenames, not independently verified
execution dates. The manifest is a retrospective selection, not a
preregistration.

\section{4. Benchmark}\label{benchmark}

The pilot covers selective evidence, scope creep, capitulation,
anchoring, sycophancy and source fabrication. Each has ten easy and ten
hard cases. Cases contain a user ask, a candidate draft, evidence
statements, an expected verdict and gold rule identifiers. The hard set
contains 32 drafts expected to be flagged and 28 expected to SHIP. Each
family's specificity denominator is only four or five clean drafts.

The cases are author-built. Near-miss clean cases share situations with
flag cases, such as disclosing a failing test rather than omitting it.
This supports inspection of over-flagging, but does not make the cases
independent samples of deployed agent failures. The release does not
provide a preregistered construction chronology establishing that
prompts and labels were fixed before case development.

External datasets such as FaithBench {[}Bao et al., 2024{]} and ELEPHANT
{[}Cheng et al., 2025{]} motivate broader factuality and sycophancy
evaluation. They are not a basis for treating these authored coding
cases as a representative sample of either benchmark. Independent
annotation and adjudication of the pilot labels are not completed study
results.

\section{5. Experiments}\label{experiments}

The main hard-set comparison has seven conditions. \texttt{none} ships
every draft without a model call. \texttt{freeform} asks for an ordinary
review. \texttt{checklist} supplies a bulk set of plain obligations.
\texttt{mono-rid} supplies the fourteen named clause families and asks
for rule IDs. \texttt{judge} checks one named target; \texttt{judge-ctx}
checks that target with the full clause list also in context.
\texttt{judge-anon} checks one plain unnamed obligation. The legacy
\texttt{contract} reviewer uses an earlier monolithic validator; it is
reported with the easy-set comparisons in the supplement, not relabelled
as the current clause-focused method.

We score exact three-way verdict accuracy, recall of gold flag cases,
specificity on gold SHIP cases and precision among predicted flags.
REVISE and BLOCK both count as flags. Attribution is reported separately
at clause-family and exact numbered-rule levels. A condition's correct
rule matches among true-positive flags use a different denominator from
its flag precision. Counts in the generated summary keep these
denominators visible.

\subsection{5.1 Detection and
specificity}\label{detection-and-specificity}

\begin{longtable}[]{@{}lrrrr@{}}
\toprule\noalign{}
Condition & Exact verdict & Recall & Specificity & Flag precision \\
\midrule\noalign{}
\endhead
\bottomrule\noalign{}
\endlastfoot
none & 28/60 & 0/32 & 28/28 & n/a \\
freeform & 40/60 & 32/32 & 16/28 & 32/44 \\
checklist & 30/60 & 32/32 & 9/28 & 32/51 \\
mono-rid & 35/60 & 32/32 & 10/28 & 32/50 \\
judge-ctx & 47/60 & 32/32 & 21/28 & 32/39 \\
judge & 47/60 & 32/32 & 23/28 & 32/37 \\
judge-anon & 48/60 & 31/32 & 24/28 & 31/35 \\
\end{longtable}

Table 2. Pooled hard-set results, 60 cases per condition. \texttt{none}
has perfect specificity because it never flags, not because it reviews
well. The unnamed control misses
\texttt{H-SCP-02\_test\_fixture\_touches\_prod}, a gold REVISE case. The
other audited conditions detect all 32 gold flags.

Because the conditions evaluate the same cases, we compare specificity
using exact paired McNemar tests on the 28 clean drafts. The named judge
exceeds \texttt{mono-rid} by 13/28, with 16 versus 3 discordant pairs (p
= 0.004). It exceeds freeform by 7/28, with 10 versus 3 discordant pairs
(p = 0.092). Wilson 95\% intervals for specificity are 0.64-0.92 for the
named judge, 0.21-0.54 for \texttt{mono-rid}, and 0.39-0.73 for
freeform; the paired tests, rather than interval overlap, assess their
differences. The freeform comparison is inconclusive at this sample
size. The unnamed-versus-named comparison has three versus two
discordant pairs (p = 1.000); this is not an equivalence test. These
exploratory p-values are unadjusted for the comparisons shown.

The context control's 21/28 specificity is consistent with a benefit
from focusing the decision, even when other clauses remain visible. It
does not isolate a universal decomposition mechanism: the targeted
reviewers know which obligation to inspect, and their instructions
differ from the bulk reviewers'. These are comparisons of known-target
review conditions, not evidence of improved general auditing.

\subsection{5.2 Rule attribution}\label{rule-attribution}

\begin{longtable}[]{@{}
  >{\raggedright\arraybackslash}p{(\columnwidth - 6\tabcolsep) * \real{0.2000}}
  >{\raggedleft\arraybackslash}p{(\columnwidth - 6\tabcolsep) * \real{0.2667}}
  >{\raggedleft\arraybackslash}p{(\columnwidth - 6\tabcolsep) * \real{0.2667}}
  >{\raggedleft\arraybackslash}p{(\columnwidth - 6\tabcolsep) * \real{0.2667}}@{}}
\toprule\noalign{}
\begin{minipage}[b]{\linewidth}\raggedright
Reviewer
\end{minipage} & \begin{minipage}[b]{\linewidth}\raggedleft
Correct family among true-positive flags
\end{minipage} & \begin{minipage}[b]{\linewidth}\raggedleft
Exact gold rule among true-positive flags
\end{minipage} & \begin{minipage}[b]{\linewidth}\raggedleft
Flag precision
\end{minipage} \\
\midrule\noalign{}
\endhead
\bottomrule\noalign{}
\endlastfoot
mono-rid & 31/32 & 28/32 & 32/50 \\
judge & 32/32 & 32/32 & 32/37 \\
\end{longtable}

Table 3. Minimal attribution checks, reported separately from false
alarms. Exact attribution requires at least one emitted numbered rule to
match a gold rule. It does not imply that every additional emitted
identifier is appropriate. The JSON supplement also reports micro
precision and recall over emitted rule sets, penalizing extra
identifiers.

Across the named judge's true-positive flags, the per-case extracted
identifier sets contain 84 identifiers, of which 37 match the 37
recorded gold identifiers. Micro precision is therefore 37/84 and micro
recall 37/37 under this literal scoring rule. The 32/32 case-level match
count does not adjudicate the appropriateness of every additional
identifier.

The distinction matters in three selective-evidence cases: the
monolithic reviewer cites \texttt{INT-SEL-02} where the gold label is
\texttt{INT-SEL-01}. A family match credits SEL; an exact-rule match
does not. The fourth monolithic attribution miss fails even the family
check. Earlier family-overlap scores such as 31/50 and 32/37 also
included false-positive flags in the denominator and should not be read
as pure rule-identification accuracy.

Rule IDs make advice referable to an external document. Comparisons with
reviewers not asked to emit IDs do not establish that naming improves
detection, nor that users benefit from the identifiers. The fair
attribution comparison is with \texttt{mono-rid}, which has the same
vocabulary.

\subsection{5.3 Severity and evidence
authority}\label{severity-and-evidence-authority}

\begin{longtable}[]{@{}lrrr@{}}
\toprule\noalign{}
Gold verdict & Predicted SHIP & Predicted REVISE & Predicted BLOCK \\
\midrule\noalign{}
\endhead
\bottomrule\noalign{}
\endlastfoot
SHIP & 23 & 4 & 1 \\
REVISE & 0 & 1 & 7 \\
BLOCK & 0 & 1 & 23 \\
\end{longtable}

Table 4. The named judge's hard-set confusion matrix. Exact accuracy is
47/60. Seven of eight gold REVISE cases are overgraded to BLOCK, and one
gold BLOCK is undergraded to REVISE. Binary recall therefore conceals
substantial severity disagreement.

The source-fabrication cases expose a second problem. Some reviewer
outputs reject supplied evidence because fixture paths are absent from
the actual host filesystem. A fluent rationale cannot resolve that
disagreement: the benchmark needs an explicit evidence authority policy,
and a deployment needs real evidence objects. The current results retain
this failure rather than treating a later policy proposal as an
implemented repair.

\section{6. Related Work}\label{related-work}

Prose-rubric review predates this contract. Constitutional AI uses
written principles in critique and training {[}Bai et al., 2022{]}.
Self-Refine {[}Madaan et al., 2023{]}, Reflexion {[}Shinn et al.,
2023{]} and CoVe {[}Dhuliawala et al., 2023{]} use feedback, reflection
or independently answered verification questions. Our reviewer belongs
to this broad family. An external document and rule-ID output are
deployment choices whose effects must be measured, not evidence of a new
factuality mechanism.

Sampling-based approaches obtain different signals. SelfCheckGPT
{[}Manakul et al., 2023{]} compares stochastic outputs, and semantic
entropy {[}Farquhar et al., 2024{]} groups answers by meaning.
Hidden-state probes estimate uncertainty from model activations
{[}Kossen et al., 2024{]}. MiniCheck instead trains a smaller
document-grounded fact checker {[}Tang et al., 2024{]}. These methods
address different access assumptions and tasks; this pilot supplies no
matched performance comparison with them. FaithBench's near-50\%
detection accuracy concerns its deliberately challenging summarization
benchmark {[}Bao et al., 2024{]}, not a universal ceiling for reviewers.

LLM-as-judge work provides the evaluation lineage {[}Liu et al., 2023;
Zheng et al., 2023{]}. Cooperative AI motivates explicit coordination
infrastructure {[}Dafoe et al., 2020{]}. Trace-Based Assurance
instruments agent orchestration and includes factuality indicators
{[}Paduraru et al., 2026{]}; RFCAudit checks protocol implementations
against RFCs {[}Zheng et al., 2025{]}; AgentVerify examines LTL safety
properties {[}Fang, 2026{]}. Their trace, specification and control-flow
units differ from our final-draft obligations. The present evidence
supports a pilot of advisory review, not a claim to supersede those
systems.

\section{7. Limitations}\label{limitations}

Oracle routing is the central limit. Each targeted reviewer receives the
case's gold mode. A deployment must select clauses or run several and
combine their outputs. Its draft-level false-positive rate, routing
errors and total cost cannot be inferred from these per-target scores.

The cases and labels were constructed within the project. Independent
human annotation, adjudication and agreement statistics such as Cohen's
kappa are required before treating those labels as externally validated.
Near-miss pairs and recurring contexts also limit independence. With
only four or five clean cases per family, one changed verdict moves a
family specificity by 0.20 or 0.25. Pooled results do not remove this
design limit.

Circularity remains: the intended drafting and auditing workflow shares
a model family, and all recorded reviewer conditions use
\texttt{sonnet}. There is no controlled cross-family replication or
repeated-run variance estimate. Missing backend and decoding metadata
limits exact reproduction even though the released outputs can be
re-scored deterministically.

Severity calibration and evidence authority remain unresolved. The 7/8
overgrade count is specific to this selected hard set, not an estimate
of operational intervention risk. Rule-ID matching does not establish
correct reasoning, and a matching ID can accompany unrelated extra IDs.
Conversely, an exact mismatch can involve similarly worded rules; the
author's gold set is not an exhaustive semantic adjudication of every
ID. The workflow has not been evaluated with human decisions or
downstream coding outcomes. It is not state of the art on factuality,
and its pilot does not validate all fourteen clauses or the catalogue's
434 records.

\section{8. Conclusion}\label{conclusion}

A written integrity contract gives a reviewer a stable obligation to
cite and gives the human an artifact to inspect and amend. In this
pilot, targeted review leaves more clean drafts unflagged than bulk
clause review, while its comparison with freeform remains inconclusive.
Naming the obligation supports exact-rule references, but the unnamed
control prevents attributing the detection pattern to catalogue names
alone. The released results support known-target review and expose
severity errors; a broader efficacy claim requires independent labels
and an end-to-end routing and aggregation evaluation.

\section{References}\label{references}

Bai, Y., Kadavath, S., Kundu, S., et al.~(2022). Constitutional AI:
Harmlessness from AI Feedback. arXiv:2212.08073v1.

Bao, F. S., Li, M., Qu, R., et al.~(2024). FaithBench: A Diverse
Hallucination Benchmark for Summarization by Modern LLMs.
arXiv:2410.13210v1.

Bradner, S. (1997). Key words for use in RFCs to Indicate Requirement
Levels. RFC 2119. https://www.rfc-editor.org/rfc/rfc2119

Cheng, M., Yu, S., Lee, C., Khadpe, P., Ibrahim, L., and Jurafsky, D.
(2025). ELEPHANT: Measuring and understanding social sycophancy in LLMs.
arXiv:2505.13995v2.

Dafoe, A., Hughes, E., Bachrach, Y., et al.~(2020). Open Problems in
Cooperative AI. arXiv:2012.08630v1.

Dhuliawala, S., Komeili, M., Xu, J., et al.~(2023).
Chain-of-Verification Reduces Hallucination in Large Language Models.
arXiv:2309.11495v2.

Fang, E. (2026). AgentVerify: Compositional Formal Verification of AI
Agent Safety Properties via LTL Model Checking. DOI:
10.20944/preprints202604.1029.v1.

Farquhar, S., Kossen, J., Kuhn, L., and Gal, Y. (2024). Detecting
hallucinations in large language models using semantic entropy. Nature
630, 625-630. DOI: 10.1038/s41586-024-07421-0.

Kossen, J., Han, J., Razzak, M., et al.~(2024). Semantic Entropy Probes:
Robust and Cheap Hallucination Detection in LLMs. arXiv:2406.15927v1.

Liu, Y., Iter, D., Xu, Y., et al.~(2023). G-Eval: NLG Evaluation using
GPT-4 with Better Human Alignment. arXiv:2303.16634v3.

Madaan, A., Tandon, N., Gupta, P., et al.~(2023). Self-Refine: Iterative
Refinement with Self-Feedback. arXiv:2303.17651v2.

Manakul, P., Liusie, A., and Gales, M. J. F. (2023). SelfCheckGPT:
Zero-Resource Black-Box Hallucination Detection for Generative Large
Language Models. arXiv:2303.08896v3.

Paduraru, C., Bouruc, P.-L., and Stefanescu, A. (2026). A Trace-Based
Assurance Framework for Agentic AI Orchestration: Contracts, Testing,
and Governance. arXiv:2603.18096v1.

Shinn, N., Cassano, F., Berman, E., et al.~(2023). Reflexion: Language
Agents with Verbal Reinforcement Learning. arXiv:2303.11366v4.

Tang, L., Laban, P., and Durrett, G. (2024). MiniCheck: Efficient
Fact-Checking of LLMs on Grounding Documents. arXiv:2404.10774v2.

Zheng, L., Chiang, W.-L., Sheng, Y., et al.~(2023). Judging
LLM-as-a-Judge with MT-Bench and Chatbot Arena. arXiv:2306.05685v4.

Zheng, M., Wang, C., Liu, X., et al.~(2025). RFCAudit: An LLM Agent for
Functional Bug Detection in Network Protocols. arXiv:2506.00714v2.

\section{Supplementary results and
reproduction}\label{supplementary-results-and-reproduction}

\begin{longtable}[]{@{}
  >{\raggedright\arraybackslash}p{(\columnwidth - 10\tabcolsep) * \real{0.1364}}
  >{\raggedright\arraybackslash}p{(\columnwidth - 10\tabcolsep) * \real{0.1364}}
  >{\raggedleft\arraybackslash}p{(\columnwidth - 10\tabcolsep) * \real{0.1818}}
  >{\raggedleft\arraybackslash}p{(\columnwidth - 10\tabcolsep) * \real{0.1818}}
  >{\raggedleft\arraybackslash}p{(\columnwidth - 10\tabcolsep) * \real{0.1818}}
  >{\raggedleft\arraybackslash}p{(\columnwidth - 10\tabcolsep) * \real{0.1818}}@{}}
\toprule\noalign{}
\begin{minipage}[b]{\linewidth}\raggedright
Suite
\end{minipage} & \begin{minipage}[b]{\linewidth}\raggedright
Reviewer
\end{minipage} & \begin{minipage}[b]{\linewidth}\raggedleft
Exact verdict
\end{minipage} & \begin{minipage}[b]{\linewidth}\raggedleft
Recall
\end{minipage} & \begin{minipage}[b]{\linewidth}\raggedleft
Specificity
\end{minipage} & \begin{minipage}[b]{\linewidth}\raggedleft
Exact rule / true-positive flags
\end{minipage} \\
\midrule\noalign{}
\endhead
\bottomrule\noalign{}
\endlastfoot
Easy & none & 24/60 & 0/36 & 24/24 & n/a \\
Easy & freeform & 44/60 & 36/36 & 19/24 & 0/36 \\
Easy & contract, legacy & 42/60 & 36/36 & 16/24 & 0/36 \\
Easy & judge & 46/60 & 36/36 & 20/24 & 34/36 \\
Hard & none & 28/60 & 0/32 & 28/28 & n/a \\
Hard & freeform & 40/60 & 32/32 & 16/28 & 0/32 \\
Hard & contract, legacy & 40/60 & 32/32 & 15/28 & 0/32 \\
Hard & judge & 47/60 & 32/32 & 23/28 & 32/32 \\
\end{longtable}

The legacy four-condition comparison is retained with exact-verdict,
recall, specificity and attribution counts for both authored suite
types. The easy set has 36 gold flags and 24 clean drafts; its label
does not imply that every condition performs better on it. Reviewers not
instructed to emit rule IDs are not a fair comparison for ID emission.

The frozen manifest and the generated summary in \texttt{experiments/}
identify each result file and condition. The summary reports the easy
suites, legacy \texttt{contract} baseline and per-family counts
alongside the hard-set ablation. The cases, raw outputs, reviewer
prompts and contract remain available for inspection. Re-scoring these
files is distinct from rerunning a model whose exact historical backend
is not recorded.

\end{document}
